\documentclass[11pt]{article}
\usepackage[margin=1in]{geometry}
\usepackage{amsmath,amssymb}
\usepackage{graphicx}
\usepackage{booktabs}
\usepackage{hyperref}
\usepackage{siunitx}
\graphicspath{{../work/}}
\title{Minimal Replication of\\
\emph{Magnetic Field Induced Incommensurate Resonance in Cuprate Superconductors}\\
\large arXiv:0805.3922v2 (Zhang, Cheng, Guo, Feng, 2008)}
\author{Automated replication (OpenClaw subagent)}
\date{2026-07-19}

\begin{document}
\maketitle

\begin{abstract}
We replicate the falsifiable analytic mechanism of arXiv:0805.3922, which studies
the dynamical spin response of cuprate superconductors under a uniform Zeeman
field within the kinetic-energy-driven superconducting (SC) mechanism (t--J model
$+$ charge--spin-separation fermion-spin theory). The paper's full self-consistent
spin self-energy (a double momentum sum with coupled gap equations) is out of
scope for a minimal analytic replication; we instead test the mechanism encoded in
its Eqs.~(4), (8), (9). We confirm the Zeeman branch splitting $\pm2\varepsilon_B$
\emph{exactly}, reproduce the field-driven commensurate$\to$incommensurate (IC)
resonance splitting and its monotonic $\delta_r(B)$ with a critical field
$\approx\SI{6.2}{\tesla}$ bracketed by the paper's $B_{c1}\approx\SI{4}{\tesla}$
and $B_{c2}\approx\SI{10}{\tesla}$, and reproduce the energy-selectivity /
hourglass-breakdown scale $\omega<0.16J$. We also report a genuine finding: the
paper's $\varepsilon_B\leftrightarrow B$ conversion implies a $g$-factor
$\approx1.04$, not the $\approx2$ expected for Cu$^{2+}$.
\textbf{Verdict: partial replication.} Coverage $7/10$, Agreement $8/10$.
\end{abstract}

\section{Note on task-label mismatch}
The task directory is named \texttt{TEXTURE-multipolar-zhang2008} and the task text
references ``multipolar texture.'' The PDF present (\texttt{paper.pdf}) is instead
arXiv:0805.3922, a cuprate spin-response paper with no multipolar-texture content.
We replicated the paper actually present and flag the mismatch (see
\texttt{extraction/marker.md}).

\section{Central claims and parameters}
The t--J model with a Zeeman term is treated in the CSS fermion-spin theory. Key
statements:
\begin{itemize}
  \item \textbf{Eq.~(4):} the mean-field spin excitation splits into two branches
  $\omega_k^{(1)}=\omega_k+2\varepsilon_B$, $\omega_k^{(2)}=\omega_k-2\varepsilon_B$,
  with $\varepsilon_B=g\mu_B B$.
  \item \textbf{Eq.~(8):} $S(k,\omega)\propto \mathrm{Im}\,\Sigma^{(s)}/
  \{[(\omega-2\varepsilon_B)^2-\omega_k^2-B_k\,\mathrm{Re}\,\Sigma^{(s)}]^2
  +[B_k\,\mathrm{Im}\,\Sigma^{(s)}]^2\}$ --- the incoming-neutron energy enters as
  $(\omega-2\varepsilon_B)$.
  \item \textbf{Eq.~(9):} resonance peaks where
  $W(k_c,\omega)=[(\omega-2\varepsilon_B)^2-\omega_{k_c}^2-B_{k_c}\mathrm{Re}\Sigma]^2\approx0$.
  \item Field induces commensurate$\to$IC resonance; $\delta_r$ increases with $B$
  (Fig.~3); critical fields $B_{c1}\approx\SI{4}{\tesla}$,
  $B_{c2}\approx\SI{10}{\tesla}$.
  \item High-energy IC scattering ($\omega\sim0.7J$) robust; hourglass breaks down
  for $\omega<0.16J\approx\SI{19}{\milli\electronvolt}$.
\end{itemize}
Parameters: $t/J=2.5$, $t'/t=0.3$, $J\approx\SI{120}{\milli\electronvolt}$,
$x=0.15$, $T=0.002J$; $\varepsilon_B=0.01J=\SI{1.2}{\milli\electronvolt}
\leftrightarrow B\approx\SI{20}{\tesla}$.

\section{Minimal model}
We use a gapped square-lattice mean-field spin mode $\omega_k$ (gap $0.31J$ at
$Q=[\pi,\pi]$, dispersing upward via $\gamma_k,\gamma'_k$), a schematic
self-energy $\Sigma$ sharply peaked at $Q$, and the exact denominator form of
Eq.~(8) including the $(\omega-2\varepsilon_B)$ shift. The
$\varepsilon_B\leftrightarrow B$ map uses the paper's own slope
$\SI{1.2}{\milli\electronvolt}/\SI{20}{\tesla}$. Full self-consistency (Eqs.~6,
7a, 7b) is out of scope.

\section{Results}
\begin{table}[h]\centering
\begin{tabular}{p{6.2cm}p{3.2cm}p{3.2cm}c}
\toprule
Claim & Paper & This work & Verdict\\\midrule
Branch splitting $=4\varepsilon_B$ (Eq.~4) & $\omega_k\pm2\varepsilon_B$ &
$0.0400$ vs $0.0400$ ($\mathrm{err}\,3\!\times\!10^{-17}$) & $\checkmark$\\
$\varepsilon_B\leftrightarrow B$ consistency & 3 anchors & single linear slope &
$\checkmark$\\
Implied $g$-factor & ($\sim2$ for Cu$^{2+}$) & $\approx1.04$ &
$\triangle$ flagged\\
Commensurate$\to$IC (raw $S$-scan) & yes & dispersion-dominated, none &
$\times$ OOS\\
Commensurate$\to$IC (Eq.~9 analytic) & $\delta_r\uparrow$ with $B$ &
commensurate $\le\SI{6}{\tesla}$, IC above, monotonic & $\checkmark$\\
Critical field & $B_{c1}4$, $B_{c2}10$~T & $B_c\approx\SI{6.2}{\tesla}$ (between) &
$\checkmark$\\
Energy selectivity $2\varepsilon_B/\omega$ & high-E robust & $2.9\%/6.5\%/20\%$
at $0.7/0.31/0.1J$ & $\checkmark$\\
Hourglass breakdown scale & $\omega<0.16J$ & $2\varepsilon_B/\omega=0.125$ at
$0.16J$ & $\checkmark$\\
\bottomrule
\end{tabular}
\caption{Quantitative comparison. OOS = out of scope (needs full self-energy).}
\end{table}

\begin{figure}[h]\centering
\includegraphics[width=0.48\textwidth]{fig_omega_spin.png}\hfill
\includegraphics[width=0.48\textwidth]{fig_delta_vs_field.png}
\caption{Left: mean-field spin excitation $\omega_k$ along $(k_x,\pi)$, gapped at
$Q$. Right: incommensurability $\delta_r$ vs field from the analytic Eq.~(9)
isolation: commensurate below $B_c\approx\SI{6.2}{\tesla}$, then monotonically
increasing IC splitting (cf.\ paper Fig.~3).}
\end{figure}

\section{Honest limitation}
The most direct check --- scanning $S(k,\omega)$ along $(k_x,\pi)$ at
$\omega=0.31J$ --- is dominated by the bare-dispersion iso-energy ring and is
essentially field-independent; it does \emph{not} reproduce the paper's
field-driven transition, because the commensurate pinning of the $B=0$ resonance
at $Q$ is a property of the self-consistent $\mathrm{Re}\,\Sigma^{(s)}$ we did not
compute. We recover the mechanism via the analytic Eq.~(9) isolation and document
the failure in \texttt{failure\_analysis.md} rather than tuning it away.

\section{Verdict}
\textbf{Partial replication: mechanism confirmed, full self-consistent map out of
scope.} \textbf{Coverage 7/10}; \textbf{Agreement 8/10}.

\end{document}
